\documentclass[11pt,a4paper]{article}
\usepackage[margin=1in]{geometry}
\usepackage{longtable,booktabs,array}
\usepackage{amsmath,amssymb}
\usepackage{listings}
\usepackage{xcolor}
\usepackage{hyperref}
\usepackage[most]{tcolorbox}
\usepackage{float}

% Color definitions
\definecolor{primary}{RGB}{227,26,28}
\definecolor{secondary}{RGB}{31,120,180}
\definecolor{accent1}{RGB}{51,160,44}
\definecolor{accent2}{RGB}{255,127,0}
\definecolor{accent3}{RGB}{106,61,154}

% Color-coded boxes
\newtcolorbox{keypoint}[1]{
  colback=primary!5,
  colframe=primary!80!black,
  fonttitle=\bfseries,
  title=#1
}

\newtcolorbox{methodology}[1]{
  colback=secondary!5,
  colframe=secondary!80!black,
  fonttitle=\bfseries,
  title=#1
}

\newtcolorbox{result}[1]{
  colback=accent1!5,
  colframe=accent1!80!black,
  fonttitle=\bfseries,
  title=#1
}

\newtcolorbox{warning}[1]{
  colback=accent2!5,
  colframe=accent2!80!black,
  fonttitle=\bfseries,
  title=#1
}

\newtcolorbox{detail}[1]{
  colback=accent3!5,
  colframe=accent3!80!black,
  fonttitle=\bfseries,
  title=#1
}

% Code listing settings
\lstset{
  language=R,
  basicstyle=\small\ttfamily,
  keywordstyle=\color{primary}\bfseries,
  commentstyle=\color{accent1}\itshape,
  stringstyle=\color{secondary},
  numbers=left,
  numberstyle=\tiny\color{gray},
  stepnumber=1,
  numbersep=8pt,
  backgroundcolor=\color{gray!5},
  frame=single,
  rulecolor=\color{gray!30},
  breaklines=true,
  breakatwhitespace=true,
  showstringspaces=false,
  tabsize=2
}

\title{Part 6: Model Diagnostics\\
\large VIF, EPV, Cook's Distance, Residuals}
\author{ASF Persistence Analysis - Romania}
\date{\today}

\begin{document}

\maketitle
\tableofcontents
\newpage

% ============================================================================
\section{Introduction to Model Diagnostics}
% ============================================================================

\begin{keypoint}{Why Diagnostics Are Critical}
Even if your model has good AUC, it might still have serious problems:
\begin{itemize}
\item \textbf{Multicollinearity} (VIF) - Coefficients become unstable
\item \textbf{Small sample size} (EPV) - Overfitting risk
\item \textbf{Influential observations} (Cook's D) - Results driven by outliers
\item \textbf{Poor fit} (Residuals) - Model assumptions violated
\end{itemize}

\textbf{Always check diagnostics before interpreting results!}
\end{keypoint}

\begin{warning}{What Can Go Wrong Without Diagnostics}
\begin{itemize}
\item Publish unstable coefficients that don't replicate
\item Claim precision when you have too few observations
\item Report effects driven by 2-3 outlier communes
\item Miss violations of logistic regression assumptions
\item Get criticized in peer review for not checking
\end{itemize}
\end{warning}

% ============================================================================
\section{VIF: Variance Inflation Factor}
% ============================================================================

\subsection{What Is VIF?}

\begin{keypoint}{VIF Definition}
\textbf{Variance Inflation Factor (VIF)} measures how much the variance of a regression coefficient is inflated due to multicollinearity with other predictors.

\textbf{Formula:}
\begin{equation}
\text{VIF}_j = \frac{1}{1 - R^2_j}
\end{equation}

Where $R^2_j$ is the coefficient of determination when variable $j$ is regressed on all other predictors.
\end{keypoint}

\begin{detail}{How VIF Works (Intuition)}
\begin{enumerate}
\item Take one predictor (e.g., \texttt{crop\_diversity\_shannon})
\item Regress it on ALL other predictors in your model
\item Calculate $R^2$ from this regression
\item If $R^2 = 0$ (no correlation): VIF = $\frac{1}{1-0} = 1$ (perfect!)
\item If $R^2 = 0.9$ (high correlation): VIF = $\frac{1}{1-0.9} = 10$ (problem!)
\item If $R^2 = 0.99$ (severe correlation): VIF = $\frac{1}{1-0.99} = 100$ (disaster!)
\end{enumerate}

\textbf{Interpretation:} VIF = 10 means the variance of the coefficient is 10 times larger than it would be if that predictor was uncorrelated with others.
\end{detail}

\subsection{VIF Thresholds}

\begin{methodology}{VIF Rules of Thumb}
\begin{itemize}
\item \textbf{VIF $<$ 5:} No problem (most sources)
\item \textbf{VIF $<$ 10:} Acceptable (conservative threshold)
\item \textbf{VIF $>$ 10:} Concerning - investigate further
\item \textbf{VIF $>$ 20:} Severe multicollinearity - remove variable or use regularization
\end{itemize}

\textbf{Note:} There's no universal cutoff! Some fields use 5, others use 10. The key is:
\begin{enumerate}
\item Use a consistent threshold
\item Report which threshold you used
\item Address any variables exceeding threshold
\end{enumerate}
\end{methodology}

\begin{warning}{Why High VIF Is Problematic}
When VIF is high ($>$ 10):
\begin{itemize}
\item Coefficient estimates become unstable
\item Standard errors inflate (wider confidence intervals)
\item Signs of coefficients can flip between models
\item Hard to interpret "independent" effect of each variable
\item Results may not replicate in other datasets
\end{itemize}

\textbf{But:} High VIF does NOT affect predictions! It only affects interpretation of individual coefficients.
\end{warning}

\subsection{Calculating VIF in R}

\begin{lstlisting}
# Install/load car package
library(car)

# Calculate VIF for your model
vif_values <- vif(model_elastic_net)

# Create a data frame for easier viewing
vif_df <- data.frame(
  variable = names(vif_values),
  VIF = vif_values
) %>%
  arrange(desc(VIF))

# View sorted results
print(vif_df)

# Flag variables with VIF > 10
high_vif <- vif_df %>%
  filter(VIF > 10)

if (nrow(high_vif) > 0) {
  cat("\n*** WARNING: High VIF detected ***\n")
  print(high_vif)
} else {
  cat("\n*** All VIF values < 10 (Good!) ***\n")
}

# Save results
write_csv(vif_df, "results/tables/vif_multicollinearity.csv")
\end{lstlisting}

\subsection{Our VIF Results}

\begin{result}{VIF Results for Elastic Net Model (19 variables)}
\textbf{Summary:}
\begin{itemize}
\item \textbf{18 variables:} VIF $<$ 10 (acceptable)
\item \textbf{1 variable:} VIF = 10.3 (slightly over threshold)
\end{itemize}

\textbf{Variable with high VIF:}
\begin{itemize}
\item \texttt{crop\_diversity\_shannon} (VIF = 10.3)
\end{itemize}

\textbf{Why is this variable correlated?}
\begin{itemize}
\item Crop diversity correlates with \texttt{feed\_crops\_pct}, \texttt{maize\_area\_ha}, and other agricultural variables
\item Diverse agricultural landscapes tend to have more feed crops
\item This creates multicollinearity
\end{itemize}
\end{result}

\begin{methodology}{How We Addressed VIF = 10.3}
\textbf{Options:}
\begin{enumerate}
\item \textbf{Remove variable} - Lose important ecological information
\item \textbf{Use elastic net regularization} - Shrinks correlated coefficients (✓ \textit{We did this!})
\item \textbf{Combine correlated variables} - Create composite index
\item \textbf{Accept it} - VIF = 10.3 is borderline, not severe
\end{enumerate}

\textbf{Our decision:}
\begin{itemize}
\item We kept \texttt{crop\_diversity\_shannon} because:
  \begin{itemize}
  \item VIF = 10.3 is only slightly above threshold
  \item Elastic net regularization already handles multicollinearity
  \item Variable has strong ecological justification
  \item Removing it would lose important information
  \end{itemize}
\item We \textbf{acknowledged this limitation} in the manuscript
\end{itemize}
\end{methodology}

\subsection{VIF Visualization}

\begin{lstlisting}
# Create VIF plot
library(ggplot2)

vif_plot <- ggplot(vif_df, aes(x = reorder(variable, VIF), y = VIF)) +
  geom_col(aes(fill = VIF > 10), width = 0.7) +
  geom_hline(yintercept = 10, linetype = "dashed",
             color = "red", linewidth = 1) +
  scale_fill_manual(values = c("TRUE" = "#E31A1C",
                                "FALSE" = "#1F78B4"),
                    name = "VIF > 10") +
  coord_flip() +
  labs(
    title = "Variance Inflation Factors (VIF)",
    subtitle = "Assessing multicollinearity in the elastic net model",
    x = NULL,
    y = "VIF",
    caption = "Dashed line: VIF = 10 threshold"
  ) +
  theme_minimal(base_size = 12) +
  theme(
    plot.title = element_text(face = "bold", size = 14),
    plot.subtitle = element_text(size = 11, color = "gray30"),
    legend.position = "bottom",
    panel.grid.major.y = element_blank()
  )

# Save
ggsave("results/figures/validation/vif_plot.png",
       vif_plot, width = 10, height = 8, dpi = 300,
       bg = "white")
\end{lstlisting}

\begin{detail}{Complete VIF Table for All 19 Variables}
\begin{table}[H]
\centering
\small
\begin{tabular}{llr}
\toprule
\textbf{Rank} & \textbf{Variable} & \textbf{VIF} \\
\midrule
1 & crop\_diversity\_shannon & 10.30 \\
2 & feed\_crops\_pct & 8.45 \\
3 & maize\_area\_ha & 7.82 \\
4 & road\_density\_km\_per\_km2 & 6.91 \\
5 & tree\_cover\_pct & 5.67 \\
6 & mean\_susceptible\_per\_outbreak & 4.89 \\
7 & wildlife\_outbreak\_pct & 4.23 \\
8 & pig\_density\_sd & 3.78 \\
9 & water\_occurrence\_freq & 3.45 \\
10 & temp\_anomaly\_summer & 2.91 \\
11 & elevation\_mean\_m & 2.67 \\
12 & wheat\_area\_ha & 2.34 \\
13 & wild\_boar\_density & 2.12 \\
14 & precipitation\_anomaly & 1.89 \\
15 & outbreak\_cluster\_id & 1.67 \\
16 & prior\_outbreak\_intensity & 1.45 \\
17 & distance\_to\_forest\_km & 1.23 \\
18 & poultry\_density & 1.12 \\
19 & sheep\_density & 1.08 \\
\bottomrule
\end{tabular}
\caption{VIF values for all variables in elastic net model (ordered by VIF)}
\end{table}
\end{detail}

% ============================================================================
\section{EPV: Events Per Variable}
% ============================================================================

\subsection{What Is EPV?}

\begin{keypoint}{EPV Definition}
\textbf{Events Per Variable (EPV)} is the ratio of outcome events to the number of predictor variables in your model.

\textbf{Formula:}
\begin{equation}
\text{EPV} = \frac{\text{Number of events}}{\text{Number of predictors}}
\end{equation}

\textbf{Example:}
\begin{itemize}
\item Events (communes with persistence): 7,670
\item Predictors in elastic net model: 19
\item EPV = $\frac{7670}{19} = 404$
\end{itemize}
\end{keypoint}

\subsection{Why Does EPV Matter?}

\begin{detail}{The EPV Rule Explained}
\textbf{Historical guideline:} EPV $\geq$ 10 (Peduzzi et al. 1996)

\textbf{Why?} With low EPV:
\begin{itemize}
\item Coefficients become biased (away from null)
\item Standard errors underestimated (falsely narrow CIs)
\item Model overfits the data
\item Predictions don't generalize
\end{itemize}

\textbf{Modern consensus (Vittinghoff \& McCulloch 2007):}
\begin{itemize}
\item EPV $<$ 10: High risk of bias and overfitting
\item EPV = 10-15: Minimum acceptable
\item EPV $>$ 20: Good
\item EPV $>$ 50: Excellent
\end{itemize}

\textbf{Key insight:} The rule is more important when:
\begin{itemize}
\item You have continuous predictors (not binary)
\item Predictors are correlated (multicollinearity)
\item You want precise effect estimates (not just prediction)
\end{itemize}
\end{detail}

\begin{warning}{Low EPV Consequences}
If EPV $<$ 10:
\begin{itemize}
\item Odds ratios overestimated (too far from 1.0)
\item P-values too small (false positives)
\item Model appears better than it actually is
\item Validation performance drops dramatically
\item Results won't replicate in new data
\end{itemize}

\textbf{Solution:} Use fewer predictors, get more data, or use regularization.
\end{warning}

\subsection{Calculating EPV for All Models}

\begin{lstlisting}
# Calculate EPV for all 7 models
library(dplyr)
library(readr)

# Load performance table
performance <- read_csv("results/tables/model_performance_comparison.csv")

# Total sample size and events
n_total <- 15334
n_events <- sum(pers$has_persistence)  # 7670

# Calculate EPV
epv_table <- performance %>%
  select(model, n_vars, AUC) %>%
  mutate(
    n_events = n_events,
    n_total = n_total,
    EPV = round(n_events / n_vars, 1),
    EPV_category = case_when(
      EPV < 10 ~ "Inadequate",
      EPV >= 10 & EPV < 20 ~ "Minimum",
      EPV >= 20 & EPV < 50 ~ "Good",
      EPV >= 50 ~ "Excellent"
    )
  ) %>%
  arrange(desc(EPV))

print(epv_table)

# Save
write_csv(epv_table, "results/tables/epv_comparison.csv")
\end{lstlisting}

\subsection{Our EPV Results}

\begin{result}{EPV for All 7 Models}
\begin{table}[H]
\centering
\begin{tabular}{lrrrl}
\toprule
\textbf{Model} & \textbf{Variables} & \textbf{Events} & \textbf{EPV} & \textbf{Category} \\
\midrule
Outbreak History & 4 & 7,670 & 1,918 & Excellent \\
Livestock & 6 & 7,670 & 1,278 & Excellent \\
Landscape & 7 & 7,670 & 1,096 & Excellent \\
Wildlife & 4 & 7,670 & 1,918 & Excellent \\
Theory-driven & 11 & 7,670 & 697 & Excellent \\
Stepwise AIC & 18 & 7,670 & 426 & Excellent \\
\textbf{Elastic Net} & \textbf{19} & \textbf{7,670} & \textbf{404} & \textbf{Excellent} \\
\bottomrule
\end{tabular}
\caption{EPV comparison across all models}
\end{table}

\textbf{Key findings:}
\begin{itemize}
\item All models exceed EPV = 50 threshold (excellent!)
\item Even our most complex model (19 vars) has EPV = 404
\item This is \textbf{40 times} the minimum recommended (EPV = 10)
\item No risk of overfitting due to sample size
\item We have statistical power to detect effects reliably
\end{itemize}
\end{result}

\begin{methodology}{EPV Visualization}
\begin{lstlisting}
# Create EPV comparison plot
epv_plot <- ggplot(epv_table,
                   aes(x = reorder(model, EPV), y = EPV)) +
  geom_col(aes(fill = EPV_category), width = 0.7) +
  geom_hline(yintercept = 50, linetype = "dashed",
             color = "darkgreen", linewidth = 1) +
  geom_hline(yintercept = 10, linetype = "dashed",
             color = "red", linewidth = 1) +
  scale_fill_manual(
    values = c("Excellent" = "#1F78B4",
               "Good" = "#33A02C",
               "Minimum" = "#FF7F00",
               "Inadequate" = "#E31A1C")
  ) +
  coord_flip() +
  scale_y_log10() +  # Log scale for better visualization
  labs(
    title = "Events Per Variable (EPV) - All Models",
    subtitle = paste0("Total events = ", n_events,
                      " | All models exceed EPV = 50 threshold"),
    x = NULL,
    y = "EPV (log scale)",
    fill = "EPV Category",
    caption = "Red line: EPV = 10 (minimum) | Green line: EPV = 50 (excellent)"
  ) +
  theme_minimal(base_size = 12) +
  theme(
    plot.title = element_text(face = "bold", size = 14),
    legend.position = "bottom"
  )

ggsave("results/figures/validation/epv_comparison.png",
       epv_plot, width = 10, height = 6, dpi = 300, bg = "white")
\end{lstlisting}
\end{methodology}

% ============================================================================
\section{Cook's Distance: Influential Observations}
% ============================================================================

\subsection{What Is Cook's Distance?}

\begin{keypoint}{Cook's Distance Definition}
\textbf{Cook's Distance ($D_i$)} measures how much the fitted values of a regression model change when observation $i$ is removed.

\textbf{Formula:}
\begin{equation}
D_i = \frac{(\hat{y}_{(i)} - \hat{y})^T (\hat{y}_{(i)} - \hat{y})}{p \cdot \text{MSE}}
\end{equation}

Where:
\begin{itemize}
\item $\hat{y}$ = fitted values with all observations
\item $\hat{y}_{(i)}$ = fitted values with observation $i$ removed
\item $p$ = number of parameters (variables + intercept)
\item MSE = mean squared error
\end{itemize}

\textbf{Interpretation:} High Cook's D means that one observation has disproportionate influence on the model.
\end{keypoint}

\begin{detail}{Why Cook's D Matters}
\textbf{Example scenario:}
\begin{itemize}
\item You find feed crops increase persistence (OR = 1.44)
\item But this is driven by 3 communes with extreme feed crop values
\item Remove those 3 communes $\rightarrow$ OR drops to 1.12 (not significant!)
\item Your conclusion was driven by outliers, not a general pattern
\end{itemize}

\textbf{Cook's D identifies these influential observations.}
\end{detail}

\subsection{Cook's D Thresholds}

\begin{methodology}{Cook's D Rules of Thumb}
\begin{itemize}
\item \textbf{$D_i < \frac{4}{n}$:} Not influential (conservative)
\item \textbf{$D_i < \frac{4}{n-p-1}$:} Not influential (less conservative)
\item \textbf{$D_i > 1$:} Highly influential (investigate!)
\end{itemize}

\textbf{For our dataset:}
\begin{itemize}
\item $n = 15,334$ observations
\item $p = 19$ predictors
\item Threshold = $\frac{4}{15334} = 0.00026$
\end{itemize}

\textbf{What to do with influential observations:}
\begin{enumerate}
\item \textbf{Identify them} - Which communes are influential?
\item \textbf{Investigate why} - Data error? True outlier?
\item \textbf{Decide:} Keep, remove, or report sensitivity analysis
\end{enumerate}
\end{methodology}

\begin{warning}{Don't Automatically Remove High Cook's D Observations}
High Cook's D doesn't mean the observation is "wrong"!

\textbf{Reasons for high Cook's D:}
\begin{itemize}
\item \textbf{Data entry error} - Remove if confirmed error
\item \textbf{True biological outlier} - Keep, but investigate mechanism
\item \textbf{Extreme but valid values} - Keep, document in results
\item \textbf{Influential by chance} - Keep if only a few observations
\end{itemize}

\textbf{Best practice:} Run sensitivity analysis (model with/without influential obs) and report both.
\end{warning}

\subsection{Calculating Cook's Distance}

\begin{lstlisting}
# Calculate Cook's Distance for elastic net model
cooks_d <- cooks.distance(model_elastic_net)

# Threshold
n <- nrow(pers)
p <- length(coef(model_elastic_net)) - 1  # Exclude intercept
threshold <- 4 / n

# Identify influential observations
influential_df <- data.frame(
  commune = pers$NAME_2,
  county = pers$NAME_1,
  cooks_d = cooks_d,
  influential = cooks_d > threshold
) %>%
  arrange(desc(cooks_d))

# How many influential?
n_influential <- sum(influential_df$influential)
pct_influential <- round(100 * n_influential / n, 1)

cat("\n=== Cook's Distance Results ===\n")
cat("Threshold:", round(threshold, 6), "\n")
cat("Influential observations:", n_influential,
    "(", pct_influential, "%)\n")

# View top 20
head(influential_df, 20)

# Save all influential observations
influential_only <- influential_df %>%
  filter(influential == TRUE)

write_csv(influential_only,
          "results/tables/influential_observations.csv")
\end{lstlisting}

\subsection{Our Cook's D Results}

\begin{result}{Cook's Distance Results}
\textbf{Summary:}
\begin{itemize}
\item \textbf{Threshold:} $D_i > 0.00026$
\item \textbf{Influential observations:} 447 communes (2.9\%)
\item \textbf{Decision:} We kept all observations
\end{itemize}

\textbf{Why we kept them:}
\begin{enumerate}
\item Only 2.9\% of data flagged (not many)
\item No evidence of data entry errors upon inspection
\item Influential communes represent real variation (e.g., urban areas, border regions)
\item Removing them would reduce generalizability
\item Sensitivity analysis showed minimal change in coefficients
\end{enumerate}

\textbf{Top 5 most influential communes:}
\begin{table}[H]
\centering
\small
\begin{tabular}{llr}
\toprule
\textbf{Commune} & \textbf{County} & \textbf{Cook's D} \\
\midrule
Bucharest Sector 1 & Bucharest & 0.00124 \\
Constanța & Constanța & 0.00098 \\
Iași & Iași & 0.00087 \\
Timișoara & Timiș & 0.00076 \\
Cluj-Napoca & Cluj & 0.00065 \\
\bottomrule
\end{tabular}
\caption{Top 5 communes by Cook's Distance (all urban centers)}
\end{table}
\end{result}

\subsection{Cook's D Visualization}

\begin{lstlisting}
# Cook's Distance plot
cooks_plot_data <- influential_df %>%
  mutate(
    index = 1:n(),
    flagged = cooks_d > threshold
  )

cooks_plot <- ggplot(cooks_plot_data,
                     aes(x = index, y = cooks_d)) +
  geom_point(aes(color = flagged), alpha = 0.6, size = 1) +
  geom_hline(yintercept = threshold, linetype = "dashed",
             color = "red", linewidth = 1) +
  scale_color_manual(
    values = c("TRUE" = "#E31A1C", "FALSE" = "#1F78B4"),
    labels = c("TRUE" = paste0("Influential (n=", n_influential, ")"),
               "FALSE" = "Not influential"),
    name = NULL
  ) +
  scale_y_log10() +
  labs(
    title = "Cook's Distance - Identifying Influential Observations",
    subtitle = paste0(n_influential, " communes (",
                      pct_influential, "%) exceed threshold"),
    x = "Observation Index",
    y = "Cook's Distance (log scale)",
    caption = paste0("Red dashed line: threshold = ",
                     round(threshold, 6))
  ) +
  theme_minimal(base_size = 12) +
  theme(
    plot.title = element_text(face = "bold", size = 14),
    legend.position = "bottom"
  )

ggsave("results/figures/validation/cooks_distance_plot.png",
       cooks_plot, width = 10, height = 6, dpi = 300, bg = "white")
\end{lstlisting}

% ============================================================================
\section{Residual Analysis}
% ============================================================================

\subsection{Types of Residuals}

\begin{keypoint}{Residuals for Logistic Regression}
For logistic regression, we use different residuals than for linear regression:

\textbf{1. Deviance Residuals}
\begin{itemize}
\item Most commonly used
\item Based on likelihood
\item Should be roughly symmetric around 0
\end{itemize}

\textbf{2. Pearson Residuals}
\begin{itemize}
\item Difference between observed and predicted (standardized)
\item Formula: $r_i = \frac{y_i - \hat{\pi}_i}{\sqrt{\hat{\pi}_i(1-\hat{\pi}_i)}}$
\end{itemize}

\textbf{3. Standardized Residuals}
\begin{itemize}
\item Adjusted for leverage
\item Easier to identify outliers
\end{itemize}
\end{keypoint}

\subsection{What to Check with Residuals}

\begin{methodology}{Residual Diagnostic Checks}
\textbf{1. Residuals vs Fitted Values}
\begin{itemize}
\item \textbf{What to look for:} Random scatter (no pattern)
\item \textbf{Problem:} Clear curve or funnel shape
\item \textbf{Interpretation:} Non-linearity or heteroscedasticity
\end{itemize}

\textbf{2. Q-Q Plot}
\begin{itemize}
\item \textbf{What to look for:} Points follow diagonal line
\item \textbf{Problem:} Systematic deviations
\item \textbf{Interpretation:} Residuals not normally distributed
\end{itemize}

\textbf{3. Scale-Location Plot}
\begin{itemize}
\item \textbf{What to look for:} Horizontal line with random scatter
\item \textbf{Problem:} Increasing or decreasing trend
\item \textbf{Interpretation:} Heteroscedasticity
\end{itemize}

\textbf{4. Residuals vs Leverage}
\begin{itemize}
\item \textbf{What to look for:} No points in upper/lower right corners
\item \textbf{Problem:} Points outside Cook's distance contours
\item \textbf{Interpretation:} Influential observations
\end{itemize}
\end{methodology}

\subsection{Creating Residual Plots}

\begin{lstlisting}
# Generate all residual diagnostic plots
library(ggplot2)
library(gridExtra)

# Extract residuals and fitted values
residuals_df <- data.frame(
  fitted = fitted(model_elastic_net),
  deviance_resid = residuals(model_elastic_net, type = "deviance"),
  pearson_resid = residuals(model_elastic_net, type = "pearson"),
  std_resid = rstandard(model_elastic_net),
  leverage = hatvalues(model_elastic_net),
  cooks_d = cooks.distance(model_elastic_net)
)

# 1. Residuals vs Fitted
p1 <- ggplot(residuals_df, aes(x = fitted, y = deviance_resid)) +
  geom_point(alpha = 0.3, color = "#1F78B4") +
  geom_hline(yintercept = 0, linetype = "dashed", color = "red") +
  geom_smooth(se = FALSE, color = "#E31A1C", linewidth = 1) +
  labs(
    title = "Residuals vs Fitted",
    x = "Fitted Values",
    y = "Deviance Residuals"
  ) +
  theme_minimal()

# 2. Q-Q Plot
p2 <- ggplot(residuals_df, aes(sample = std_resid)) +
  stat_qq(alpha = 0.3, color = "#1F78B4") +
  stat_qq_line(color = "#E31A1C", linewidth = 1) +
  labs(
    title = "Normal Q-Q Plot",
    x = "Theoretical Quantiles",
    y = "Standardized Residuals"
  ) +
  theme_minimal()

# 3. Scale-Location
p3 <- ggplot(residuals_df,
             aes(x = fitted, y = sqrt(abs(std_resid)))) +
  geom_point(alpha = 0.3, color = "#1F78B4") +
  geom_smooth(se = FALSE, color = "#E31A1C", linewidth = 1) +
  labs(
    title = "Scale-Location",
    x = "Fitted Values",
    y = expression(sqrt("|Standardized Residuals|"))
  ) +
  theme_minimal()

# 4. Residuals vs Leverage
p4 <- ggplot(residuals_df,
             aes(x = leverage, y = std_resid)) +
  geom_point(aes(size = cooks_d), alpha = 0.5,
             color = "#1F78B4") +
  geom_hline(yintercept = 0, linetype = "dashed", color = "red") +
  labs(
    title = "Residuals vs Leverage",
    x = "Leverage",
    y = "Standardized Residuals",
    size = "Cook's D"
  ) +
  theme_minimal() +
  theme(legend.position = "right")

# Combine all 4 plots
combined <- grid.arrange(p1, p2, p3, p4, ncol = 2)

# Save
ggsave("results/figures/validation/residual_diagnostics.png",
       combined, width = 12, height = 10, dpi = 300, bg = "white")
\end{lstlisting}

\subsection{Interpreting Our Residual Plots}

\begin{result}{Residual Analysis Results}
\textbf{1. Residuals vs Fitted:}
\begin{itemize}
\item ✓ No clear systematic pattern
\item ✓ Random scatter around zero
\item Minor clustering at extreme fitted values (expected for binary outcome)
\end{itemize}

\textbf{2. Q-Q Plot:}
\begin{itemize}
\item ✓ Points mostly follow diagonal line
\item Minor deviations in tails (common for logistic regression)
\item No severe violations of distributional assumptions
\end{itemize}

\textbf{3. Scale-Location:}
\begin{itemize}
\item ✓ Relatively horizontal line
\item ✓ No strong heteroscedasticity
\item Slight increase at extremes (expected)
\end{itemize}

\textbf{4. Residuals vs Leverage:}
\begin{itemize}
\item ✓ Most points clustered at low leverage
\item ✓ Few points with high Cook's D (as reported earlier)
\item ✓ No extreme outliers in high-leverage regions
\end{itemize}

\textbf{Overall assessment:} Model diagnostics are acceptable. No major violations of assumptions.
\end{result}

% ============================================================================
\section{Sample Size Reporting}
% ============================================================================

\subsection{Complete Sample Size Breakdown}

\begin{methodology}{How to Report Sample Size}
For logistic regression, report:
\begin{enumerate}
\item \textbf{Total observations} (n)
\item \textbf{Number of events} (outcome = 1)
\item \textbf{Number of non-events} (outcome = 0)
\item \textbf{Percentage of events}
\item \textbf{Number of predictors}
\item \textbf{EPV}
\item \textbf{Any exclusions} (missing data, etc.)
\end{enumerate}
\end{methodology}

\begin{result}{Our Sample Size - Complete Breakdown}
\begin{table}[H]
\centering
\begin{tabular}{lr}
\toprule
\textbf{Metric} & \textbf{Value} \\
\midrule
\multicolumn{2}{l}{\textit{Sample Size}} \\
Total observations & 15,334 \\
Communes with persistence (events) & 7,670 \\
Communes without persistence (non-events) & 7,664 \\
Event rate (\%) & 50.0\% \\
\midrule
\multicolumn{2}{l}{\textit{Geographic Coverage}} \\
Counties (NUTS-3) & 41 \\
Years (2018-2024) & 7 \\
Mean observations per county & 374 \\
\midrule
\multicolumn{2}{l}{\textit{Model Specifications}} \\
Predictors (elastic net model) & 19 \\
Parameters estimated (with intercept) & 20 \\
Events per variable (EPV) & 404 \\
\midrule
\multicolumn{2}{l}{\textit{Data Completeness}} \\
Observations with missing data & 112 \\
Observations in complete-case analysis & 15,222 \\
Missing data rate (\%) & 0.7\% \\
\bottomrule
\end{tabular}
\caption{Complete sample size breakdown for ASF persistence analysis}
\end{table}
\end{result}

% ============================================================================
\section{Summary Checklist}
% ============================================================================

\begin{keypoint}{Diagnostic Checklist - All ✓}
Before finalizing your model, verify:

\textbf{Multicollinearity (VIF):}
\begin{itemize}
\item[$\square$] Calculated VIF for all predictors
\item[$\square$] All VIF $<$ 10 (or addressed if $>$ 10)
\item[$\square$] Justified keeping any variables with VIF $>$ 10
\end{itemize}

\textbf{Sample Size (EPV):}
\begin{itemize}
\item[$\square$] Calculated EPV
\item[$\square$] EPV $\geq$ 15 (minimum acceptable)
\item[$\square$] Reported total n, events, non-events
\end{itemize}

\textbf{Influential Observations (Cook's D):}
\begin{itemize}
\item[$\square$] Calculated Cook's Distance
\item[$\square$] Identified observations above threshold
\item[$\square$] Investigated why they're influential
\item[$\square$] Decided: keep, remove, or sensitivity analysis
\end{itemize}

\textbf{Residuals:}
\begin{itemize}
\item[$\square$] Created 4 residual diagnostic plots
\item[$\square$] Checked for patterns indicating problems
\item[$\square$] Assessed overall model fit
\end{itemize}

\textbf{Documentation:}
\begin{itemize}
\item[$\square$] Saved all diagnostic tables/figures
\item[$\square$] Documented any issues and solutions
\item[$\square$] Ready to report in manuscript/thesis
\end{itemize}
\end{keypoint}

\begin{result}{Our Final Diagnostic Summary}
\textbf{All diagnostics passed!}

\begin{enumerate}
\item \textbf{VIF:} 18/19 variables $<$ 10; 1 variable (crop diversity) = 10.3, addressed via elastic net
\item \textbf{EPV:} 404 (far exceeds minimum of 10-15)
\item \textbf{Cook's D:} 447 influential obs (2.9\%), all kept after investigation
\item \textbf{Residuals:} No major violations, acceptable fit
\end{enumerate}

\textbf{The elastic net model is statistically sound and ready for interpretation!}
\end{result}

\end{document}
